\begin{statement}{C.2.UN.22}{Proved}{A component cover with a rank-two translation frame has coherent Euler characteristic zero}
Let $X$ be a smooth compact complex surface with
$\kappa(X)\ge0$. Suppose that it carries a two-dimensional
complex vector space $\mathfrak a$ of commuting meromorphic
vector fields, of generic rank two, every member of which
is semicomplete on its holomorphic locus. Then
\[
                         \chi(\mathcal O_X)=0.
                                                               \tag{C.RF.5}
\]
More precisely, $X$ has a two-torus birational model,
or it has an actual relatively minimal elliptic fibration
on a proper birational model, over a smooth compact curve
of genus at most one, with identity integral $H_1$ monodromy
and with every fibre a positive multiple of a smooth elliptic
curve. Its smooth fibres are torsors under one fixed elliptic
curve. All actual multiple fibres and their multiplicities
are retained in this conclusion.

This applies in particular to every covering component
$E$ of C.2.UN.16 with $\kappa(E)\ge0$ in its
generic-rank-two branch. Such a covering component cannot
have a K3 or Enriques minimal model. The assertion concerns
the actual specified root-cover component; quotients of that
component are compared separately.
\end{statement}

\begin{proof}[The elliptic-first-integral step in C.2.UN.22]
Let $X$ be a smooth compact complex surface and let $V,Y$ be
commuting meromorphic vector fields with $V\wedge Y\ne0$.
Assume that both fields are semicomplete on their holomorphic
loci and that an actual resolution of a meromorphic first
integral of $V$, followed by its Stein factorization, gives
a map
\[
 q:\widehat X\longrightarrow B
 \tag{C.RF.1}
\]
with connected elliptic general fibre.  The proof below does
not require $X$ to be K\"ahler.  We use the meromorphic lifts
of the two fields on $\widehat X$; their brackets still vanish
on the common open isomorphism and hence everywhere, and
their semicompleteness extends over their holomorphic loci
by \cite[Theorem~3.3]{Guillot2024}.

First, $V$ has no horizontal polar curve.  On a general
smooth fibre $F$ of $q$, its restriction is a nonzero
semicomplete meromorphic field on the compact elliptic
curve.  The time-form argument in C.2.UN.16 excludes poles
on $F$.  Thus this restriction is holomorphic and nowhere
zero, and its flow acts by all translations of $F$.
The commutation relation implies
$(\Phi_V^a)_*Y=Y$ as an identity of meromorphic fields.
Consequently the intersection of the polar divisor of $Y$
with $F$ is invariant under all translations of $F$.
It is a finite set, so it is empty.  This proves that
$Y$ also has no horizontal polar curve.

The field $Y$ projects to a nonzero meromorphic field $Z$
on $B$.  In fact, for each local meromorphic function $g$
on $B$,
\[
 V\bigl(Y(q^*g)\bigr)
   =Y\bigl(V(q^*g)\bigr)=0.
\]
The meromorphic function $Y(q^*g)$ is therefore constant
on the connected general fibres; proper Stein descent
makes it the pullback of a meromorphic function on $B$.
These functions define $Z$.  Were $Z$ zero, both fields
would be tangent to the one-dimensional fibres, contrary
to $V\wedge Y\ne0$.

The projected field has no pole.  To verify this without
an unproved descent assertion about semicompleteness,
suppose that $Z$ has a pole of order $r\ge1$ at $b$.
The primitive of its time form has a zero of order $r+1$
there.  Its inverse solution consequently has distinct
continuations after a loop about the corresponding time
value.  Choose that loop and its lifted base paths in a
small punctured disc about $b$ avoiding all other bad
values.  The inverse image of the compact union of these
paths is compact.  Both $q$ and $Y$ are holomorphic there,
and $Y$ has no pole there, since all its polar curves are
vertical and the disc was chosen accordingly.  The usual
continuation theorem for holomorphic ordinary differential
equations therefore lifts the whole time loop to a
solution of $Y$.  Its endpoint has a different image in
$B$.  This contradicts semicompleteness of $Y$.
Hence $Z$ is a nonzero holomorphic field on the compact
curve $B$, and $g(B)\le1$.

Remove finitely many values so that $q$ is a smooth
elliptic fibration, $V,Y$ are holomorphic, $V$ is nowhere
zero vertically, and $Z$ is nonzero.  On this open base
$B^\circ$, let $\alpha_b$ be the time form of $V|_{F_b}$.
The flow of $Y$, defined for a uniform small time along
each compact fibre, carries whole fibres biholomorphically
to whole fibres and carries $V$ to itself.  It therefore
preserves their time forms.  Integration gives the exact
horizontal map from the integral homology local system
\[
 \mathcal H_1(q)\big|_{B^\circ}
   \longrightarrow\underline{\mathbb C},
 \qquad
 \gamma\longmapsto\int_\gamma\alpha_b.
 \tag{C.RF.2}
\]
The map is injective on each stalk: the periods of a
nonzero holomorphic one-form on an elliptic curve are a
rank-two lattice in $\mathbb C$.  The preceding flow
comparison proves that its values are locally constant.
Analytic continuation of any integral cycle around a
loop preserves its complex period; injectivity then
proves that every monodromy operator on $H_1(F_b,\mathbb Z)$
is the identity.  The image lattice is fixed, say
$\Lambda\subset\mathbb C$, and the fibres over
$B^\circ$ are torsors under the fixed elliptic curve
$\mathbb C/\Lambda$.

Contract the vertical exceptional curves of the first
kind to obtain the actual relatively minimal elliptic
map
\[
 \widehat X\xrightarrow{\ \mu\ }X_0
       \xrightarrow{\ q_0\ }B.
 \tag{C.RF.3}
\]
These contractions preserve the preceding integral
monodromy representation.  Kodaira's classification
\cite[Theorem~6.2, pp.~564--566, and Table~I, p.~604]{Kodaira1963II}
now shows that every fibre has smooth elliptic support.
Here is the comparison for multiplicities, to specify
the use of the table.  The only multiple types in that
classification are $mI_b$.  Locally write such a full
fibre as $t=u h^m$, where $h=0$ is its reduced Cartier
support and $u$ is a unit.  In the normalization of the
actual base change $t=s^m$, the charts
$w^m=u$, $s=wh$ are \emph{\'etale} over the original
smooth surface.  The covering family has a reduced
central fibre and its linear monodromy is the $m$th
power of the original one, hence the identity.  It
remains relatively minimal: every central component
downstairs has canonical degree zero by the displayed
Kodaira types; the \emph{\'etale} canonical pullback
gives the same zero degree on every central prime
curve upstairs, and adjunction excludes an exceptional
rational curve of self-intersection $-1$ upstairs.
Among the reduced relatively minimal types in the
cited table, identity monodromy is the smooth type
$I_0$.  Hence the original support is smooth as well.
The argument applies to every connected covering
component and makes no connectedness assumption on
the root cover.

In particular, for pairwise distinct base values the
fibre supports $C_b$ are disjoint smooth elliptic
curves, and $C_b^2=0$, since some positive multiple
of $C_b$ is a full fibre.  The Euler characteristic
is zero by additivity over the smooth base stratum
and the finitely many special fibres.  To compute
the canonical square, choose a nonzero meromorphic
one-form $\eta$ on $B$.  The meromorphic field $V$
descends through $\mu$, and at the generic point
there is a unique meromorphic two-form $\Omega$
satisfying
\[
 \iota_V\Omega=q_0^*\eta.
 \tag{C.RF.4}
\]
Both sides are interpreted as actual meromorphic
differentials.  Since $V$ is tangent to $q_0$, the
right-hand side annihilates $V$, which proves
existence and uniqueness of $\Omega$ by a local
tangent-frame calculation.  Its divisor is
vertical: on a general compact fibre $V$ is
holomorphic and nowhere zero, and $\eta$ is
holomorphic and nonzero at its base value.
Thus a canonical divisor is a finite integral
combination of the disjoint square-zero curves
$C_b$, and $K_{X_0}^2=0$.  The surface
Riemann--Roch formula gives
$\chi(\mathcal O_{X_0})=(K_{X_0}^2+e(X_0))/12=0$.
Proper birational invariance of the coherent
Euler characteristic under the displayed point
blowups and blowdowns yields
$\chi(\mathcal O_X)=0$.

No vanishing of $Z$ at a multiple value is asserted.
Indeed, a logarithmic elliptic translation can absorb
the cyclic root monodromy.  Multiple smooth fibres
are included in the conclusion.
\end{proof}

\begin{proof}[Completion of C.2.UN.22]
We first treat the nonprojective cases by actual
meromorphic maps. If $a(X)=1$, choose a nonconstant
$h\in\mathcal M(X)$, resolve its graph, and take the
Stein factor $q:\widehat X\to B$. Every
$g\in\mathcal M(X)$ is algebraic over $\mathbb C(h)$.
On a connected general fibre of $q$ it therefore takes
finitely many values and is constant. Proper meromorphic
descent consequently gives the exact field identity
\[
                       \mathcal M(X)=q^*\mathcal M(B).
                                                               \tag{C.RF.9}
\]
For example, this descent can be seen by mapping the
proper graph of $g$ to $B\times\mathbb P^1$; its image
is the graph of the descended meromorphic function.
Each field in $\mathfrak a$ differentiates this field
of functions and therefore induces a meromorphic vector
field on $B$. The induced fields commute. They cannot
both vanish, since that would make both original fields
tangent to the one-dimensional fibres and violate
generic rank two. On a curve any two commuting fields
with one nonzero are constant multiples, by (C.RF.8).
A nonzero constant linear combination in $\mathfrak a$
is thus vertical. On the compact general fibre its
semicompleteness excludes poles, so the genus is at most
one; genus zero again contradicts $\kappa(X)\ge0$ by
adjunction. This gives an elliptic fibration and reduces
the case $a(X)=1$ to (C.RF.1)--(C.RF.4).

If $a(X)=0$, apply the final argument of C.2.UN.16.
That argument uses only the given two-dimensional
space of semicomplete fields: the fibration alternative
of \cite[Theorem~9.2]{Guillot2024} is impossible in
algebraic dimension zero, and the fields on their
respective holomorphic birational models descend to
two independent holomorphic fields on the same unique
nonruled minimal model. The compact-surface field
classification and dimension bounds used there give
a two-torus or a Hopf surface. Kodaira dimension
excludes the latter. This gives the two-torus alternative
and (C.RF.5).

It remains to treat $a(X)=2$. A compact complex surface
of algebraic dimension two has a smooth projective
birational model. Passing through the common resolution
preserves the bracket, generic rank, and coherent Euler
characteristic; semicompleteness on the full holomorphic
locus follows from the same closedness comparison as in
C.2.UN.16. We may therefore work on that projective model.
Choose a basis $V_1,V_2$ of $\mathfrak a$. The nonzero
meromorphic field $V_1$ has either a leaf with Zariski-dense
image or only algebraic leaves. In the latter case the
constructible integral-curve incidence argument in
C.2.UN.15 gives an actual resolution of its first integral
and a Stein fibration. The semicomplete time-form calculation
on its compact general fibre gives genus at most one.
A rational general fibre would force every pluricanonical
section to vanish, by adjunction as in C.2.UN.20, contrary
to $\kappa(X)\ge0$. Thus its general fibre is elliptic,
and the preceding proof applies to $V_1,V_2$.

It remains to treat a Zariski-dense leaf of $V_1$.
Semicompleteness supplies the univalent maximal solution
required by \cite[Main Theorem and the clarification immediately
following it, p.~2]{Guillot2019}. On an actual smooth
birational model, either the surface is Abelian, when
(C.RF.5) follows at once, or there is a fibration
\[
                         q:X'\longrightarrow B
                                                               \tag{C.RF.6}
\]
over a rational or elliptic curve, with rational or elliptic
general fibre, which $V_1$ preserves and whose fibres
contain every polar curve of $V_1$. Again the rational-fibre
case contradicts $\kappa(X')\ge0$. Thus (C.RF.6) is an
elliptic fibration. The projection $Z_1=q_*V_1$ is nonzero:
otherwise every leaf of $V_1$ would lie in a compact fibre,
contrary to the selected dense leaf. Both lifted fields
remain semicomplete on their holomorphic loci by the same
birational-open comparison and closedness argument used
in C.2.UN.16.

We prove that the polar divisor of $V_2$ is also vertical.
At a generic point where $V_1$ is holomorphic and nonzero,
the identity $[V_1,V_2]=0$ says that its local flow preserves
$V_2$ as a meromorphic tensor and therefore preserves its
polar divisor. Equivalently, in a flow coordinate
$V_1=\partial_z$, the coefficients of $V_2$ are independent
of $z$. Thus a horizontal polar component would be a
compact invariant multisection of the foliation of $V_1$.
We now rule out such a multisection with the actual finite
maps needed for this conclusion.

Shrink to a nonempty open base $B^\circ$ on which the
fibration is smooth, $V_1$ is holomorphic on each whole
fibre, and $Z_1$ is nonzero. Its flow is defined uniformly
for small time on any one compact fibre, and maps that
whole fibre biholomorphically to a neighboring fibre.
Consequently the elliptic curves are isotrivial over
$B^\circ$. Suppose that $C\subset X'$ were a compact
horizontal invariant curve. Its normalization $\widetilde C$
maps finitely to $B$. In the normalized fibre product with
$\widetilde C$, the diagonal curve supplies an actual
section; it is horizontal for the lifted field. Resolve
this fibre product, taking a resolution which is an
isomorphism above the smooth open locus.

Fix an elliptic curve $A$ isomorphic to a smooth fibre.
Over a smaller open subset of $\widetilde C$, the
origin-preserving isomorphisms from $A$ to the fibres
form a finite \'etale torsor with structure group
$\operatorname{Aut}(A,0)$. This is a concrete finite
cover: in a Weierstrass equation with constant $j$, the
change of coordinates is obtained by adjoining the
scalars $\lambda$ with
$g_2=\lambda^4g_{2,A}$ and
$g_3=\lambda^6g_{3,A}$, omitting the equation with
zero constant coefficient when $j=0$ or $1728$.
Its ambiguity is exactly the finite origin-preserving
automorphism group. Normalize the compact curve
$\widetilde C$ in a connected component of this finite
cover. On the resulting compact curve $B'$ there is
an actual product isomorphism over a nonempty open set,
and resolving its graph gives proper maps
\[
 \begin{array}{ccc}
 \widetilde X'&\longrightarrow&X'\\
 \downarrow&&\downarrow q\\
 B'&\longrightarrow&B,
 \end{array}
 \qquad
 \widetilde X'\dashrightarrow A\times B'
 \quad\hbox{birational}.
                                                               \tag{C.RF.7}
\]
Here all curves and surfaces are projective, so the
isomorphism over the algebraic open set defines the
displayed rational map; its graph closure is proper.
The section is carried to the zero section.

On the product open set the lifted field has the form
$a(b)\partial_b+c(b)\partial_u$, where $\partial_u$
is a fixed translation field on $A$. Indeed, its vertical
coefficient is holomorphic on each whole compact elliptic
fibre, so is constant on that fibre. Tangency to the zero
section gives $c=0$. Every local leaf therefore lies in
$\{u\}\times B'$. The compact closure of this section
in the proper graph model maps to an algebraic invariant
curve in $X'$. Through a general point this contains
the entire leaf by uniqueness of analytic continuation.
This contradicts the dense leaf of $V_1$. This argument
also explains the invariant-multisection step in
\cite[\S6.3]{Guillot2019}; no extension of an arbitrary
transcendental horizontal trivialization is being used.

It follows that $V_2$ has no horizontal polar curve.
Hence it is holomorphic on every fibre over a suitable
nonempty open subset of $B^\circ$. Its normal component
there is a holomorphic function on a compact elliptic
curve and is constant. It consequently projects to a
meromorphic field $Z_2$ on $B$; the coefficients descend
through the proper connected-fibre map, or can be
recovered by the trace along any finite multisection.
The bracket identity gives $[Z_1,Z_2]=0$. For nonzero
$Z_1$ on a connected curve this is the precise equation
\[
              Z_1(Z_2/Z_1)=0,
              \qquad Z_2=cZ_1\quad(c\in\mathbb C).
                                                               \tag{C.RF.8}
\]
Thus $W=V_2-cV_1$ is vertical. It is nonzero and
$W\wedge V_1\ne0$ by the generic-rank-two hypothesis.
It is semicomplete because it is a constant member of
$\mathfrak a$, and it commutes with $V_1$. Apply the
preceding elliptic-first-integral proof to $W,V_1$
and (C.RF.6). It gives the stated relatively minimal
model, fixed period lattice, identity integral monodromy,
and (C.RF.5).

Finally $\chi(\mathcal O)=2$ on a K3 surface and
$\chi(\mathcal O)=1$ on an Enriques surface. These
values and coherent Euler characteristic are preserved
under the actual point blowups and blowdowns comparing
smooth birational surface models. They exclude those
two minimal models for the rank-two covering component.
\end{proof}


% C.2.RAT.1 from part_C2_rational.tex

